\begin{helper}
Let \(D\) be a digraph on a finite vertex set \(W\), let \(r:W\to\mathbb N\) be a ranking function, and let \(S\subseteq W\) be an independent \(k\)-set in the ordered graph \(\Gamma(D,r)\) from \(\noderef{DigraphOrderedGraph}\). Suppose that \(r\) is injective on \(S\). Then there is an injective ordered tuple
\[
  v:\{1,\ldots,k\}\to W
\]
whose range is exactly \(S\) and which is forward independent in \(D\), in the sense of \(\noderef{ForwardIndependentTuple}\).
\end{helper}
\begin{proof}
Since \(S\) is an independent \(k\)-set in \(\Gamma(D,r)\), it is a \(k\)-clique in the complement of \(\Gamma(D,r)\). Thus \(S\) has exactly \(k\) elements, and every two distinct vertices of \(S\) are nonadjacent in \(\Gamma(D,r)\).

Because \(r\) is injective on \(S\), the image \(r(S)\subseteq\mathbb N\) also has exactly \(k\) elements. List the elements of \(r(S)\) in increasing order as
\[
  \rho_1<\rho_2<\cdots<\rho_k .
\]
For each \(i\), choose the unique vertex \(v_i\in S\) with \(r(v_i)=\rho_i\). The uniqueness follows from the injectivity of \(r\) on \(S\). The tuple \(v_1,\ldots,v_k\) is injective, and its range is exactly \(S\): every \(v_i\) lies in \(S\), and if \(x\in S\), then \(r(x)\in r(S)\), so \(r(x)=\rho_i\) for some \(i\), and injectivity on \(S\) gives \(x=v_i\).

It remains to check forward independence. Let \(i<j\). Then \(r(v_i)=\rho_i<\rho_j=r(v_j)\). If \((v_i,v_j)\in A(D)\), then the definition of \(\Gamma(D,r)\) in \(\noderef{DigraphOrderedGraph}\) puts an edge between \(v_i\) and \(v_j\), because the lower-ranked vertex \(v_i\) sends an arc to the higher-ranked vertex \(v_j\). This contradicts the fact that the two distinct vertices \(v_i,v_j\in S\) are nonadjacent in \(\Gamma(D,r)\). Therefore no such arc \((v_i,v_j)\) exists for any \(i<j\), which is exactly the condition from \(\noderef{ForwardIndependentTuple}\).
\end{proof}
